\newglossaryentry{term-onchain-data}{
    name={Ончейн-данные},
    description={данные, зафиксированные в блокчейне, а также производные представления, формируемые на их основе для прикладного анализа}
}
\newglossaryentry{term-indexer}{
    name={Индексер},
    description={серверное приложение, которое получает данные блокчейна, нормализует их, сохраняет в базе данных и предоставляет прикладной API}
}
\newglossaryentry{term-bitcoin-core}{
    name={Bitcoin Core},
    description={программная реализация узла сети Bitcoin, используемая как доверенный источник первичных блоков и транзакций}
}
\newglossaryentry{term-bitcoin}{
    name={Bitcoin},
    description={децентрализованная блокчейн-сеть и одноименная криптовалютная система, данные которой являются предметом индексации}
}
\newglossaryentry{term-canonical}{
    name={Каноническая цепь},
    description={актуальная ветвь блокчейна Bitcoin, признаваемая основной в рассматриваемый момент времени}
}
\newglossaryentry{term-job}{
    name={Задание индексирования},
    description={управляемый процесс обработки блокчейн-данных с собственным идентификатором, статусом, прогрессом и параметрами выполнения}
}
\newglossaryentry{term-reorg}{
    name={Reorg},
    description={смена канонической ветви блокчейна, из-за которой ранее принятые блоки и транзакции становятся неканоническими}
}
\newglossaryentry{term-mempool}{
    name={Mempool},
    description={множество неподтвержденных транзакций, известных узлу Bitcoin Core и ожидающих включения в блок}
}
\newglossaryentry{term-stateless}{
    name={Stateless-сервис},
    description={сервис, который не хранит критичное состояние локально между перезапусками и размещает его во внешнем хранилище}
}
\newglossaryentry{term-address-index}{
    name={Адресный индекс},
    description={структура данных, позволяющая быстро находить транзакции, выходы, балансы и mempool-события, связанные с конкретным адресом}
}
\newglossaryentry{term-address-balance}{
    name={Адресный баланс},
    description={агрегированное значение суммы непотраченных выходов или исторического состояния адреса на определенный момент цепи}
}
\newglossaryentry{term-utxo}{
    name={UTXO},
    description={непотраченный выход транзакции, который может быть использован как вход последующей транзакции}
}
\newglossaryentry{term-postgresql}{
    name={PostgreSQL},
    description={реляционная система управления базами данных, используемая для хранения блоков, транзакций, UTXO, балансов, заданий и метрик состояния}
}
\newglossaryentry{term-rust}{
    name={Rust},
    description={язык программирования, на котором реализована серверная часть системы}
}
\newglossaryentry{term-python}{
    name={Python},
    description={язык программирования, используемый для реализации командного интерфейса проекта}
}
\newglossaryentry{term-react}{
    name={React},
    description={JavaScript-библиотека, используемая при реализации административной панели}
}
\newglossaryentry{term-vite}{
    name={Vite},
    description={инструмент сборки frontend-приложений, применяемый для административной панели}
}
\newglossaryentry{term-docker}{
    name={Docker},
    description={платформа контейнеризации, применяемая для воспроизводимой сборки и развертывания компонентов системы}
}
\newglossaryentry{term-docker-compose}{
    name={Docker Compose},
    description={средство описания и запуска набора связанных контейнеров проекта}
}
\newglossaryentry{term-nginx}{
    name={nginx},
    description={веб-сервер и reverse proxy, используемый в проектной схеме защищенного контура и контейнерного развертывания}
}
\newglossaryentry{term-prometheus}{
    name={Prometheus},
    description={система мониторинга, для которой приложение публикует метрики состояния индексатора и его фоновых процессов}
}
\newglossaryentry{term-satoshi}{
    name={Satoshi},
    description={минимальная неделимая единица Bitcoin, используемая для представления значений выходов и балансов}
}
\newglossaryentry{term-openapi}{
    name={OpenAPI},
    description={спецификация описания HTTP API, используемая для документирования контрактов backend-сервиса}
}
\newglossaryentry{term-swagger-ui}{
    name={Swagger UI},
    description={веб-интерфейс для просмотра и интерактивного вызова API на основе OpenAPI-описания}
}
\newglossaryentry{term-regtest}{
    name={Regtest},
    description={локальный режим сети Bitcoin, предназначенный для воспроизводимых тестов с управляемой генерацией блоков}
}
\newglossaryentry{term-testcontainers}{
    name={Testcontainers},
    description={инструмент запуска временных контейнеров в автоматизированных тестах, используемый для проверки сценариев с PostgreSQL}
}
